\begin{textsample}{POS Dim 1 – vem\_ed\_26 – Score 25.00 – vem\_ed\_26/t138.txt}  \label{ex:f1_pos_060}
IVEC 2024: \textbf{pontes} virtuais A conferência de abertura da IVEC 2024, principal evento internacional de Intercâmbios Virtuais, \textbf{foi} proferida em 21 de outubro pela consultora e educadora feminista Keshia Abraham.
Nesta edição, a IVEC \textbf{foi} totalmente on-line.
Intitulada “Under the Bridge and Above the Muck: Engaging Virtual Exchange Holistically”, a palestra \textbf{foi} repleta de imagens potentes, como a da “\textbf{ponte} sobre o pântano”.
Muck faz \textbf{referência} às Everglades, região pantanosa no sul da Flórida, chamada “rio de grama” pela escritora e ambientalista Marjorie Stoneman Douglas.
Os Intercâmbios Virtuais \textbf{são} descritos, geralmente, como \textbf{forma} de \textbf{construir} \textbf{pontes}.
Mas não se pode erguer \textbf{pontes} sobre pântanos. \textbf{É} preciso aprofundar e conhecer o terreno primeiro, ou \textbf{seja}, o \textbf{contexto} dos aprendizes – olhar por debaixo das \textbf{pontes}, na \textbf{proposta} de Keshia Abraham. “Sobre o quê estamos \textbf{construindo}?
Com qual intenção?”, interrogou a conferencista. “Precisamos olhar sob as \textbf{pontes}, com uma curiosidade intencional e um envolvimento com o \textbf{que} \textbf{cria} essas \textbf{pontes}: como \textbf{são} feitas, quem está \textbf{construindo}, onde estão localizadas, quem tem quais recursos: eles \textbf{são} equitativos?”, questionou.
Como passos para \textbf{construir} a \textbf{ponte}, ela sugeriu: Planejamento Fundação: implantar os alicerces do \textbf{programa} Instalar os pilares e suportes (redes de \textbf{apoio} para estudantes e parceiros internacionais) Completar a \textbf{estrutura} (\textbf{implementar} o Intercâmbio Virtual de \textbf{forma} \textbf{que} \textbf{todos} possam participar de \textbf{forma} equitativa) Verificar a qualidade e segurança do projeto \textbf{continuação} Ao \textbf{longo} dos quatro dias de evento, a equipe dos PCIs fez cinco participações na IVEC 2024: Osvaldo Succi Junior mediou duas discussões facilitadas, uma delas com Regiane Camargo e outra com as parceiras Eva Haug, Rosi León e Ntsoaki Malebo; Patrícia Patrício participou de mesa-redonda sobre Inteligência Artificial nos Intercâmbios Virtuais e teve trabalho selecionado para"Flash Session"assíncrona; Ana Carolina Freschi e Neusa Gritti tiveram trabalho selecionado para"Flash Session"assíncrona.
Em 21 de outubro, Patrícia Patrício apresentou o trabalho desenvolvido em coautoria com Ana Torres e Paula Pinheiro (Universidade de Aveiro, Portugal) e Carlos Menezes (Fatec Ipiranga, Brasil).
O título da apresentação \textbf{foi}"AI-Enhanced Perception Analysis in COIL: Unveiling the Dynamics of a Long-standing Partnership"(Análise de percepção aprimorada por IA em projetos COIL: Revelando a dinâmica de uma parceria de \textbf{longa} data).
Trata-se de uma síntese dos principais resultados obtidos pelo uso de ferramentas de IA para análise das Pesquisas de Percepção respondidas pelos alunos da Fatec Ipiranga e da Universidade de Aveiro \textbf{que} participaram do PCI / Cesu realizado entre as duas Instituições de Ensino Superior no segundo semestre de 2023.
A análise exploratória \textbf{foi} realizada por estudantes do curso de Big Data para Negócios da Fatec Ipiranga e, sob orientação do professor Menezes, buscou compreender aspectos como gerenciamento do tempo, liderança, consciência intercultural, comunicação, trabalho em equipe e aprendizagem.
Utilizando a ferramenta DoWhy Library, a análise mostrou \textbf{que} \textbf{é} \textbf{possível} inferir \textbf{que} trabalho em equipe e liderança influenciam 50 \% da aprendizagem.
Consciência intercultural e liderança têm o mesmo peso.
No dia 22 de outubro, Osvaldo Succi Junior (coordenador dos PCIs / Cesu / CPS), Ntsoaki Malebo (Central University of Technology, África do Sul), Eva Haug (Amsterdam University of Applied Sciences, Países Baixos) e Rosi León (DePaul University, EUA) abordaram o \textbf{desenvolvimento} \textbf{profissional} para projetos COIL (Collaborative Online International Learning) sob uma \textbf{perspectiva} \textbf{comparada} entre o Norte e o Sul Global.
Os apresentadores analisaram as diferentes modalidades de \textbf{desenvolvimento} \textbf{profissional} adotadas por suas instituições em uma tabela comparativa, levando em conta seus \textbf{pontos} fortes e desafios \textbf{específicos}.
Depois disso, \textbf{pequenos} grupos com professores do Norte e Sul Global refletiram sobre suas práticas, considerando suas particularidades e desafios.
Osvaldo Succi Junior sintetizou: “de maneira \textbf{geral}, o Norte \textbf{é} mais planejado e o Sul, mais flexível”.
Em 23 de outubro, Osvaldo Succi Junior e Regiane Camargo conduziram a discussão sobre os \textbf{papéis} da Inteligência Artificial nos projetos de Intercâmbio Virtual, focando em três aspectos principais: o equilíbrio das participações de falantes nativos e não nativos nas \textbf{interações} escritas; o auxílio aos Intercâmbios Virtuais / COIL e a identificação das limitações da IA. “Assessment and Evaluation in Virtual Exchange Projects: a perspective of Centro Paula Souza practice” (Avaliação em projetos de \textbf{intercâmbio} virtual: uma \textbf{perspectiva} da prática do Centro Paula Souza) \textbf{é} o título da apresentação assíncrona preparada por Ana Carolina Freschi e Neusa Gritti.
O trabalho apontou exemplos de práticas de avaliação dos PCIs, como feedback (de professores para alunos, e entre pares) e a realização das Pesquisas de Percepção semestrais com estudantes e professores para análise detalhada dos PCIs / Cesu.
O instrumento de Pesquisa de Percepção está \textbf{sendo} reformulado para agregar mais \textbf{perspectivas} teóricas. \textbf{continuação} As professoras Lilian de Souza e Vera Jordan, da Fatec Bragança Paulista, experientes em PCIs / Cesu, também tiveram trabalho selecionado para “Flash Session” assíncrona, intitulado “PCI entre Brasil e Chile: uma conversa em 3 \textbf{idiomas}”. \textbf{É} o resumo de um PCI trilíngue (espanhol, português e inglês) aliando professores de línguas estrangeiras e de disciplinas técnicas do curso de Análise e Desenvolvimento de Sistemas da Fatec Bragança Paulista.
O projeto \textbf{foi} realizado em colaboração com o professor Felipe Zambrano e seus alunos da DUOC UC (Chile), abordando o tema da economia circular com \textbf{foco} em \textbf{desenvolvimento} \textbf{sustentável}.
Um relato sobre esse projeto \textbf{foi} \textbf{publicado} nas páginas 56 e 57 do boletim 66 “Conectados al Territorio”, da Duoc UC, em maio de 2024: https://observatorio.duoc.cl/boletin/boletin-n66-conectados-al-territorio-experiencias-de-vcm-en-duoc-uc/

% matched lemmas: apoio, comparar, construir, contexto, continuação, criar, desenvolvimento, específico, estrutura, foco, forma, geral, idioma, implementar, interação, intercâmbio, longo, papel, pequeno, perspectiva, ponte, ponto, possível, profissional, programa, proposta, publicar, que, referência, ser, sustentável, todo
\end{textsample}
